\FloatBarrier
\section{مقایسه‌ی نقشه‌های \lr{rollout} با پژوهش‌های پیشین}
\label{sec:literature}

بخش~\ref{sec:consistency} نشان داد کدام بخش‌های نقشه‌ها بین شرکت‌کننده‌ها تکرار می‌شوند. اما تکرارپذیر بودن یعنی پایدار بودن، نه درست بودن. این بخش می‌پرسد: آیا همین بخش‌های تکرارپذیر با آنچه پژوهش‌های قبلی درباره‌ی رابطه‌ی \lr{EEG} و فعالیت عضله یافته‌اند می‌خوانند؟

\paragraph{دو دسته پژوهش مرجع.}
\begin{itemize}
\item \textbf{\lr{corticomuscular coherence (CMC)}.} این روش اصلی ساختن نقشه‌ی «از عضله به مغز» است. \lr{CMC} می‌سنجد نوسان‌های یک کانال \lr{EEG} (یا \lr{MEG}) و نوسان‌های \lr{EMG} یک عضله در یک فرکانس خاص چقدر با هم هماهنگ‌اند. اگر یک ناحیه‌ی قشر و یک عضله هم‌آهنگ نوسان کنند، احتمالاً آن ناحیه فرمان عضله را می‌فرستد.
\item \textbf{\lr{event-related desynchronization (ERD)}.} این پژوهش‌ها تغییر توان باندهای \lr{EEG} را در اطراف حرکت می‌سنجند. مثلاً توان باند بتا روی قشر حرکتی کمی پیش از شروع حرکت پایین می‌آید.
\end{itemize}

\paragraph{یک تفاوت مهم در روش.} مدل \lr{v2} توان باندهای \lr{EEG} را می‌بیند، یعنی اینکه نوسان در هر باند چقدر بزرگ است (با هموارسازی ۵ هرتز)، نه فاز آن. \lr{CMC} اما به هم‌فازی نوسان‌های قشر و عضله در ۱۵ تا ۳۰ هرتز بستگی دارد. هدف مدل هم پوش \lr{EMG} با \lr{low-pass} ۸ هرتز است که این نوسان‌های سریع را ندارد. پس مدل اصلاً نمی‌تواند \lr{CMC} را یاد بگیرد و آنچه یاد می‌گیرد به \lr{ERD} نزدیک‌تر است. در نتیجه مقایسه با \lr{CMC} فقط درباره‌ی این است که اطلاعات «کجا» است، نه اینکه «چه» اطلاعاتی است.

\subsection{مکان: کدام الکترودها}

\paragraph{پژوهش‌های پیشین.} \lr{CMC} در باند بتا بیشترین مقدار را روی قشر حسی-حرکتی نیمکره‌ی مقابل دست دارد. برای گرفتن دقیق با دست راست، \cite{chen2013} \lr{CMC} معنادار را روی \lr{M1} نیمکره‌ی چپ و روی ناحیه‌ی حرکتی مکمل (\lr{SMA}) گزارش کرده‌اند. روی همین داده (\lr{WAY-EEG-GAL})، \cite{cisotto2018} \lr{CMC} را بین \lr{C3} و هر پنج عضله حساب کرده‌اند؛ البته فقط برای یک شرکت‌کننده (\lr{P7}) و با \lr{C3} که از پیش انتخاب شده بود. در پژوهشی با ۳۳ نوع گرفتن \cite{sburlea2018}، \lr{EEG} ناحیه‌ی آهیانه‌ای نیمکره‌ی مقابل در باند \lr{mu} در مرحله‌های پایانی گرفتن و نگه‌داشتن با فعالیت عضله رابطه داشت.

\paragraph{نتیجه‌ی ما.} تنها الکترودهایی که در سطح گروه به‌طور معنادار سهمی بیشتر از یکنواخت دارند \lr{C3} و \lr{CP1} هستند (جدول~\ref{tab:cons-top}). این دو دقیقاً روی قشر حسی-حرکتی و حسی-آهیانه‌ای نیمکره‌ی چپ‌اند، یعنی نیمکره‌ی مقابل دست راست. این با هر سه گزارش بالا سازگار است. \lr{C4} در نیمکره‌ی هم‌سمت دست سوم است ولی معنادار نیست؛ پژوهش‌های \lr{CMC} هم سازمان دوطرفه با برتری نیمکره‌ی مقابل دیده‌اند. ناحیه‌ی \lr{SMA} (اطراف \lr{Cz} و \lr{FCz}) در نقشه‌ی ما برجسته نیست.

\paragraph{آنچه با پژوهش‌ها نمی‌خواند.} در بسیاری از افراد، الکترودهای پس‌سری (\lr{O1}، \lr{PO9}، \lr{PO10})، گیجگاهی و پیشانی هم سهم بالایی دارند. این نواحی در پژوهش‌های \lr{CMC} دست جایی ندارند. به احتمال زیاد این سهم از پردازش بینایی (دیدن چراغ شروع و جسم)، حرکت چشم یا فعالیت عضلات سر و گردن می‌آید، نه از فرمان حرکتی.

\subsection{فرکانس: کدام باندها}

\paragraph{پژوهش‌های پیشین.} \lr{CMC} در باند بتا (۱۵ تا ۳۰ هرتز) هنگام انقباض پایدار و نگه‌داشتن بیشترین است و با شروع حرکت کم می‌شود \cite{kilner2000}. روی همین داده، طیف \lr{CMC} بین \lr{C3} و عضله‌ها قله‌هایی در حدود ۱۰، ۲۰ و ۳۰ هرتز داشت \cite{cisotto2018}. در \lr{ERD}، توان \lr{mu} و بتا روی قشر حسی-حرکتی مقابل از حدود یک ثانیه پیش از حرکت کم می‌شود، و پس از پایان حرکت توان بتا حتی از حالت پایه بالاتر می‌رود \cite{pfurtscheller1999}.

\paragraph{نتیجه‌ی ما.} در \lr{gradient × input}، باند بتا (۱۳ تا ۳۰ هرتز) با میانگین \lr{۳۳٫۵} درصد بیشترین سهم را دارد و در ۹ نفر از ۱۲ نفر اول است. در ۳ نفر دیگر آلفا (۸ تا ۱۳ هرتز) اول است (جدول~\ref{tab:cons-bands}). این ترتیب با هر دو دسته پژوهش سازگار است. سهم گامای پایین (۳۰ تا ۴۰ هرتز، حدود ۲۵ درصد) بیشتر از انتظار است. این باند به \lr{artifact} عضلات سر حساس است.

\subsection{زمان: پیش یا پس از رخداد}

\paragraph{پژوهش‌های پیشین.} \lr{ERD} بتا حدود یک ثانیه پیش از شروع حرکت آغاز می‌شود \cite{pfurtscheller1999}. \lr{CMC} برعکس، در نگه‌داشتن پایدار بیشترین و در حین حرکت کمترین است \cite{kilner2000}.

\paragraph{نتیجه‌ی ما.} قله‌ی \lr{gradient-weighted rollout} خروجی \lr{E} به‌طور میانگین \lr{۰٫۵} تا \lr{۰٫۸} ثانیه پیش از اولین تماس و \lr{lift-off} است (جدول~\ref{tab:cons-time}). این با زمان‌بندی \lr{ERD} پیش از حرکت می‌خواند. اما دو نکته با \lr{CMC} نمی‌خواند. اول، زمان این قله بین افراد خیلی جابه‌جا می‌شود (انحراف معیار \lr{۱٫۲} تا \lr{۱٫۶} ثانیه). دوم، \lr{E} در دو ثانیه‌ی پس از \lr{lift-off} کمترین سهم را به \lr{tokens} می‌دهد؛ درست همان آغاز نگه‌داشتن که \lr{CMC} در پژوهش‌ها بیشترین است. پس مدل بیشتر از تغییرهای اطراف رویدادهای حرکتی استفاده می‌کند تا از هماهنگی پایدار مغز و عضله در نگه‌داشتن.

\subsection{عضله: کدام عضله‌ها}

\paragraph{پژوهش‌های پیشین.} \lr{CMC} برای عضله‌های انتهایی دست (مثل \lr{FDI}) قوی و شناخته‌شده است. روی همین داده، بهترین طبقه‌بندی وزن جسم با \lr{CMC} از \lr{BR}، \lr{FDI} و \lr{FD} به دست آمد و \lr{AD} ضعیف‌ترین بود \cite{cisotto2018}.

\paragraph{نتیجه‌ی ما.} ترتیب برعکس است. \lr{E} برای \lr{AD} (عضله‌ی شانه، نزدیک تنه) بهترین پیش‌بینی را دارد ($r$ میانگین \lr{۰٫۷۸}) و برای عضله‌های انگشت، \lr{FD} و \lr{FDI}، ضعیف‌ترین ($r$ حدود \lr{۰٫۵۰}). نقشه‌های الکترود پنج عضله هم تقریباً یکسان‌اند (جدول~\ref{tab:cons-elec-musc}).

این دو نتیجه با هم نشان می‌دهند که \lr{E} بیشتر مرحله‌ی کلی حرکت بازو را از \lr{EEG} می‌خواند (دست در حال رفتن است، جسم را گرفته، بلندش کرده یا در استراحت است)، نه فرمان قشری خاص هر عضله را. احتمالاً \lr{AD} که بازو را بالا می‌برد بیشترین تغییر را بین این مرحله‌ها دارد و برای همین با این اطلاعات بهتر از همه پیش‌بینی می‌شود. همین توضیح می‌دهد چرا \lr{EEG} تقریباً چیزی به \lr{K} اضافه نمی‌کند (بخش~\ref{sec:residual}): مرحله‌ی حرکت در سینماتیک هم هست.

\subsection{تفاوت بین افراد}

پژوهش‌های \lr{CMC} تفاوت زیادی بین افراد گزارش می‌کنند. حتی در جوانان سالم، اندازه‌ی \lr{CMC} از فردی به فرد دیگر بسیار فرق می‌کند و در برخی افراد اصلاً معنادار نیست \cite{ushiyama2011,matsuya2017}. نتیجه‌ی ما هم همین تصویر را دارد: نقشه‌ی هر فرد فقط کمی شبیه دیگران است (\lr{pairwise $\rho$} حدود \lr{۰٫۲})، ولی نقشه‌ی گروه تکرارپذیر است (\lr{split-half} حدود \lr{۰٫۷}). پس ادعاهای مکانی باید درباره‌ی گروه باشند.

\subsection{مقایسه با \lr{Pinto Neto (2026)} روی همین داده}

نزدیک‌ترین پژوهش به کار ما \cite{pintoneto2026} است. آن پژوهش روی همین ۱۲ شرکت‌کننده‌ی \lr{WAY-EEG-GAL} ده مدل را آزموده است، از جمله یک \lr{graph attention network} روی گرافی از ۳۲ کانال \lr{EEG} و ۵ کانال \lr{EMG} و یک \lr{Transformer} چهارلایه. تفسیر مدل‌ها هم با \lr{feature importance}، فیلترهای \lr{EEGNet} و \lr{attention rollout} انجام شده است.

مسئله‌ی آن پژوهش با ما فرق دارد. آن‌ها از یک پنجره‌ی یک‌ثانیه‌ای در زمان نگه‌داشتن، وزن و جنس سطح جسم را طبقه‌بندی می‌کنند و مدل را روی هر نفر با مدلی که روی ۱۱ نفر دیگر آموزش دیده می‌سنجند (\lr{leave-one-subject-out}). ما پوش پیوسته‌ی \lr{EMG} را برای هر نفر جداگانه از سینماتیک و \lr{EEG} پیش‌بینی می‌کنیم. با این حال چند نتیجه‌ی آن‌ها مستقیم با نتایج ما قابل مقایسه است (جدول~\ref{tab:pinto}).

\begin{table}[htbp]
\centering\footnotesize
\caption{مقایسه با \lr{Pinto Neto (2026)} روی \lr{WAY-EEG-GAL}.}
\label{tab:pinto}
\begin{tabular}{p{0.15\textwidth}p{0.38\textwidth}p{0.38\textwidth}}
\toprule
جنبه & \lr{Pinto Neto (2026)} & این گزارش \\
\midrule
سهم \lr{EEG} & \lr{EEG} به‌تنهایی نزدیک شانس (دقت \lr{۰٫۳۴} تا \lr{۰٫۳۸}، شانس \lr{۰٫۳۳})؛ افزودن \lr{EEG} به \lr{EMG} دقت را کم می‌کند & \lr{F} برابر \lr{K}؛ \lr{EEG} حدود \lr{۰٫۵} درصد واریانس باقی‌مانده را توضیح می‌دهد (بخش~\ref{sec:residual}) \\
کنار گذاشتن \lr{EEG} & \lr{graph attention} با کم کردن وزن کانال‌های کم‌اطلاعات در برابر افزودن \lr{EEG} مقاوم است & \lr{gate} نزدیک صفر می‌ماند؛ همان سازوکار به شکل صریح \\
\lr{EEG} به‌تنهایی & در نگه‌داشتن پایدار، وزن و سطح از \lr{EEG} قابل خواندن نیست & \lr{E} با $r$ \lr{۰٫۶۴۵} مرحله‌ی حرکت را در طول کل سری می‌خواند \\
الکترودها & بیشترین \lr{importance} \lr{EEG} برای سطح روی \lr{FC1}، \lr{Cz}، \lr{C3}، \lr{FC2} و \lr{C4}؛ به گفته‌ی خودشان نشانه‌ی هم‌تغییری با \lr{EMG}، نه رمزگشایی مستقل & کانون گروهی \lr{C3} و \lr{CP1}، سپس \lr{C4}؛ همان نوار حسی-حرکتی \\
فرکانس & فیلترهای \lr{EEGNet} در بتا و گامای بالا؛ برترین ویژگی‌های \lr{EEG} گامای بالا در \lr{Fp1}، \lr{FC5} و \lr{O2} & بتا اول؛ گامای پایین حدود ۲۵ درصد؛ سهم پیشانی و پس‌سری \\
عضله & برترین ویژگی وزن: توان گامای بالای \lr{EMG} عضله‌ی \lr{AD} & \lr{E} بهترین برای \lr{AD} \\
\lr{attention rollout} & توجه دریافتی گره‌های \lr{EMG} بیشتر از گره‌های \lr{EEG}؛ میانگین روی همه‌ی \lr{fold}ها، بدون سنجش یکسانی بین افراد & سنجش یکسانی بین ۱۲ نفر (بخش~\ref{sec:consistency}) \\
\bottomrule
\end{tabular}
\end{table}

\paragraph{کجا هم‌نظریم.} دو پژوهش با مدل‌ها و مسئله‌های متفاوت به یک نتیجه می‌رسند: در \lr{WAY-EEG-GAL} اطلاعات اصلی در سیگنال‌های بدن است (\lr{EMG} در آن پژوهش، سینماتیک در ما)، و \lr{EEG} سطح سر چیز قابل‌اندازه‌گیری‌ای به آن اضافه نمی‌کند. رفتار مدل‌ها هم شبیه است: \lr{graph attention} آن‌ها به کانال‌های \lr{EEG} توجه کمی می‌کند و \lr{gate} ما نزدیک صفر می‌ماند. نقشه‌ی الکترود هر دو کار روی نوار حسی-حرکتی است و \lr{C3} و \lr{C4} در هر دو دیده می‌شوند.

\paragraph{کجا فرق داریم.} \lr{E} ما با $r$ \lr{۰٫۶۴۵} از هر ۱۲ \lr{null test} گذشت، اما \lr{EEG} آن پژوهش نزدیک شانس بود. این تناقض نیست. پنجره‌ی آن‌ها فقط زمان نگه‌داشتن است و می‌پرسند جسم سنگین است یا سبک. در این بازه دست ثابت است و تفاوت شرایط بسیار کوچک است. \lr{E} ما در طول کل سری تفاوت مرحله‌های حرکت را می‌خواند (رسیدن، گرفتن، بلند کردن، استراحت) که تفاوت‌های بسیار بزرگ‌تری‌اند. این با نتیجه‌ی بخش‌های قبل هم می‌خواند: \lr{E} بیشتر مرحله‌ی کلی حرکت را از \lr{EEG} می‌گیرد تا فرمان خاص عضله را.

\paragraph{چه چیزی اضافه می‌کنیم.} \lr{Pinto Neto} نقشه‌های توجه را فقط توصیفی و به‌صورت میانگین روی همه‌ی \lr{fold}ها گزارش کرده و سنجش پایداری آن‌ها بین افراد را به کار آینده سپرده است. بخش~\ref{sec:consistency} دقیقاً همین را برای نقشه‌های ما انجام می‌دهد: نقشه‌ی گروه تکرارپذیر است (\lr{split-half} حدود \lr{۰٫۷}) و نقشه‌ی یک نفر پرنویز.

یک تفاوت در پیش‌پردازش هم هست. آن پژوهش عمداً \lr{ICA} انجام نداده و \lr{EMG} را یکسوسازی نکرده است؛ در خط پایه‌ی ما \lr{ICA} انجام شده و هدف پوش هموار است. پس سهم بالای گامای پیشانی و پس‌سری در آن پژوهش ممکن است بیشتر از نتایج ما از \lr{artifact} اثر گرفته باشد.

\subsection{جمع‌بندی و محدودیت‌ها}

\begin{table}[htbp]
\centering\footnotesize
\caption{خلاصه‌ی مقایسه‌ی نتایج \lr{rollout} با پژوهش‌های پیشین.}
\label{tab:lit}
\begin{tabular}{p{0.15\textwidth}p{0.34\textwidth}p{0.3\textwidth}p{0.11\textwidth}}
\toprule
جنبه & پژوهش‌های پیشین & نتیجه‌ی ما & سازگاری \\
\midrule
مکان & قشر حسی-حرکتی مقابل (\lr{M1})، \lr{SMA}، آهیانه‌ای مقابل & \lr{C3} و \lr{CP1} معنادار؛ \lr{C4} سوم؛ \lr{SMA} برجسته نیست & بله، با یک استثنا \\
\lr{artifact} & -- & سهم پس‌سری، گیجگاهی و پیشانی در بسیاری از افراد & هشدار \\
فرکانس & بتا (و \lr{mu}) & بتا اول در ۹ نفر از ۱۲، آلفا در ۳ نفر & بله \\
زمان & \lr{ERD} از حدود ۱ ثانیه پیش از حرکت؛ \lr{CMC} بیشینه در نگه‌داشتن & قله حدود \lr{۰٫۵} تا \lr{۰٫۸} ثانیه پیش از رخداد؛ کمترین سهم در آغاز نگه‌داشتن & با \lr{ERD} بله، با \lr{CMC} نه \\
عضله & \lr{CMC} قوی در عضله‌های انتهایی دست & \lr{E} بهترین برای \lr{AD}، ضعیف برای \lr{FD} و \lr{FDI}؛ نقشه‌ی خاص هر عضله نیست & نه \\
تفاوت افراد & زیاد & متوسط در سطح فرد، تکرارپذیر در سطح گروه & بله \\
\bottomrule
\end{tabular}
\end{table}

جدول~\ref{tab:lit} نتیجه را خلاصه می‌کند. نقشه‌های \lr{rollout} در مکان و فرکانس با دانش موجود سازگارند: کانون تکرارپذیرشان روی قشر حسی-حرکتی مقابل دست است و باند بتا را ترجیح می‌دهند. در زمان و در خاص بودن برای هر عضله با \lr{CMC} نمی‌خوانند، و به \lr{ERD} و خواندن مرحله‌ی کلی حرکت نزدیک‌ترند.

این مقایسه کیفی است، نه عددی. هیچ‌کدام از این پژوهش‌ها نقشه‌ی عددی‌ای برای \lr{WAY-EEG-GAL} با ۱۲ شرکت‌کننده منتشر نکرده‌اند که بتوان مستقیم با آن همبستگی گرفت. پژوهش \lr{CMC} روی همین داده \cite{cisotto2018} فقط یک نفر و یک الکترود را بررسی کرده است. نقشه‌های \lr{Pinto Neto} \cite{pintoneto2026} هم برای مسئله‌ی دیگری (طبقه‌بندی وزن و سطح در نگه‌داشتن) و فقط به‌صورت میانگین گروه گزارش شده‌اند.

گام بعدی منطقی ساختن نقشه‌های مرجع روی خود داده‌ی ماست: \lr{CMC} (از \lr{EMG} خام، چون پوش هموار باند بتا ندارد) و \lr{ERD}، برای همین ۱۲ شرکت‌کننده و همین رخدادها. سپس می‌توان همبستگی مکانی این نقشه‌ها با نقشه‌های \lr{rollout} را برای هر فرد اندازه گرفت. این کار در بخش \ref{sec:controls-ref} انجام شده است.
